\chapter{Schizophrenia}
\index{Schizophrenia}
\label{sec:Schizophrenia}

\section{Overview}
Schizophrenia has a lifetime prevalence of 0.3--0.7\%, with a male-to-female ratio of 1.4:1, with men typically presenting earlier (peak onset 20--28 years) than women (peak onset 25--35 years). This condition is among the most common mental health disorders worldwide. It affects people across all age groups, socioeconomic backgrounds, and geographic regions. Mental health conditions affect thoughts, emotions, behaviors, and overall well-being. These conditions are among the most common health concerns worldwide, affecting people across all age groups. Mental health disorders result from complex interactions of biological, psychological, and social factors. Effective treatments are available, and recovery is possible with appropriate support.
\section{Causes}
This condition happens because of dopaminergic hyperactivity in the mesolimbic pathway (positive symptoms), dopaminergic hypoactivity in the mesocortical pathway (negative symptoms and thinking and memory deficits), glutamatergic NMDA receptor hypofunction, and widespread gray matter loss, with genetic factors accounting for 80\% of risk.  Neurotransmitter imbalances (serotonin, dopamine, norepinephrine) play key roles in many mental health conditions. Genetic factors contribute to vulnerability, with heritability estimates varying by condition. Early life experiences, trauma, and chronic stress can alter brain development and function. Environmental factors including social support, socioeconomic status, and life events influence mental health.   
\section{Risk Factors}
\begin{itemize}[noitemsep]
  \item Genetic predisposition and family history
  \item Advancing age
  \item Lifestyle factors (diet, exercise, smoking, alcohol)
  \item Environmental exposures
  \item Underlying medical conditions
  \item Medication use
\end{itemize}
\section{Signs and Symptoms}
\subsection{Common symptoms}
\begin{itemize}[noitemsep]
  \item \item You may experience two or more of the following for a significant portion of time during a 1-month period (at least one must be delusions, hallucinations, or disorganized speech): delusions (fixed false beliefs, often bizarre), hallucinations (auditory most common), disorganized speech (loose associations, tangentiality), grossly disorganized or catatonic behavior, and negative symptoms (blunted affect, avolition, alogia, anhedonia, asociality), with continuous signs persisting for at least 6 months.
  \item Pain or discomfort in the affected area
  \end{itemize}
\subsection{Emergency warning signs}
\begin{itemize}[noitemsep]
  \item Severe or worsening symptoms of schizophrenia require immediate medical evaluation
\end{itemize}
\section{Progression and Stages}
The condition may progress through different stages of severity over time. Early stages often involve mild or subtle symptoms that may be overlooked. Moderate stages involve more noticeable symptoms affecting daily function. Advanced stages may involve significant impairment and complications.
\section{Complications}
\begin{itemize}[noitemsep]
  \item Disease progression leading to worsening symptoms
  \item Secondary infections or injuries
  \item Side effects of treatment
  \item Reduced quality of life
  \item Emotional and psychological impact
\end{itemize}
\section{Diagnosis}
Doctors usually diagnose this based on your symptoms and a physical exam---no biomarkers are available.

\begin{itemize}[noitemsep]
  \item Comprehensive medical history and physical examination
  \item Laboratory tests to assess organ function
  \item Imaging studies to visualize affected structures
  \item Specialized testing depending on the affected system
  \item Consultation with relevant specialists
\end{itemize}
\section{Prevention}
\begin{itemize}[noitemsep]
  \item Maintain a healthy lifestyle with balanced diet and exercise
  \item Avoid known risk factors
  \item Attend regular medical check-ups for early detection
  \item Follow recommended screening guidelines
\end{itemize}
\section{Treatment Options}
Treatment options include antipsychotics (both first-generation and second-generation, with clozapine being the only FDA-approved medication for treatment-resistant schizophrenia), thinking and memory-behavioral therapy for psychosis (CBTp), psychosocial interventions, and thinking and memory remediation therapy. Psychotherapy (cognitive behavioral therapy, interpersonal therapy, psychodynamic therapy) Pharmacotherapy (antidepressants, antipsychotics, mood stabilizers, anxiolytics) Combined treatment approaches for optimal outcomes Hospitalization for acute crisis management Support groups and peer support programs Lifestyle interventions (exercise, sleep hygiene, stress management) Mindfulness and meditation practices Family therapy and psychoeducation Vocational and social rehabilitation Long-term maintenance therapy for chronic conditions

\begin{itemize}[noitemsep]
  \item Medications to manage symptoms and address underlying mechanisms
  \item Lifestyle modifications and self-care measures
  \item Physical therapy and rehabilitation
  \item Surgical intervention when indicated
  \item Regular monitoring and follow-up care
\end{itemize}
\section{Living with It}
\begin{itemize}[noitemsep]
  \item Follow your treatment plan as prescribed
  \item Maintain regular follow-up with your healthcare provider
  \item Make necessary lifestyle adjustments
  \item Seek support from family, friends, or support groups
  \item Stay informed about your condition
\end{itemize}
\section{Outlook}
The outlook varies from person to person: 20--30\% have a good outcome, with a chronic course and functional decline in many, and suicide risk of 5--10\%. With appropriate treatment, most people with mental health conditions experience significant improvement. Early intervention is associated with better long-term outcomes. Mental health conditions are chronic for some individuals, requiring ongoing management. Reducing stigma and improving access to care remain important public health priorities.
